
\section{Client On-Boarding}\label{sec:on-boarding}

\textit{Normative language per \href{https://tools.ietf.org/html/rfc2119}{RFC~2119}.}

\subsection{Overview}

Onboarding is a one-time, out-of-band process that MUST be completed before any tenant
user attempts to authenticate against Canvasflow. It consists of three phases executed
in sequence. Phases~1 and~2 may proceed in parallel for distinct workstreams, but
Phase~3 (validation) cannot begin until both are complete.

The Canvasflow OAuth Validator (\url{https://oauth.canvasflow.work}) is the mandatory
gate for production go-live. No tenant may transition to \texttt{status = 'active'}
without producing all-green results on validation checks V1 through V5.

\subsection{Phase~1 --- Canvasflow-Side Setup}\label{sec:canvasflow-side-setup}

Canvasflow completes the following steps before the tenant begins Phase~2.



\textbf{Step~1 --- Assign resource identifiers.} For every publication or content
resource the tenant is authorized to access, Canvasflow generates a stable Canvasflow
resource identifier using the \texttt{cf:<type>:<id>} format (e.g.,
\texttt{cf:issue:1111}). These identifiers are permanent. They MUST NOT change after
onboarding without prior written notice to the tenant's designated technical lead
(\hyperref[sec:canvasflow-guarantees]{\S11 Canvasflow Guarantees~6}). Once assigned, every identifier is visible in 
the Canvasflow administration dashboard — the tenant's 
technical lead retrieves them directly from the dashboard UI during Phase 2 (Step 2.1).

\textbf{Step~2 --- Create the tenant configuration record.} Canvasflow creates a row in
\texttt{tenant\_config} with:
\begin{itemize}
  \item \texttt{iss}: the tenant's IdP issuer URL (obtained from the tenant in advance)
  \item \texttt{signing\_type}: provisionally set to \texttt{jwks}; updated to
    \texttt{hs256} if the tenant cannot support JWKS (explicit exception, requires
    Canvasflow approval)
  \item \texttt{client\_id}: the OAuth client ID for the Canvasflow SPA in the tenant's
    IdP
  \item \texttt{audience}: \texttt{canvasflow-api}
  \item \texttt{status}: \texttt{pending} --- the tenant cannot authenticate against
    production until this is changed to \texttt{active} at the conclusion of Phase~3
\end{itemize}

Signing credentials (\texttt{jwks\_url} or \texttt{hs256\_secret\_encrypted}) are left
empty until the tenant provides them in Step~2.5 of Phase~2.

\textbf{Step~3 --- Provide sandbox credentials.} Canvasflow provides the tenant's
technical lead with:
\begin{itemize}
  \item The Canvasflow sandbox GraphQL endpoint URL
  \item A test user account in the Canvasflow sandbox environment for end-to-end
    validation
\end{itemize}

\subsection{Phase~2 --- Tenant-Side Setup}\label{sec:tenant-side-setup}

The tenant's technical lead performs the following steps inside their own IdP. The order
is sequential --- each step depends on the previous.

\textbf{Step~2.1 --- Retrieve resource identifiers from the Canvasflow administration dashboard.}

Log into the Canvasflow administration dashboard and navigate to the publications 
or content resources the tenant is authorized to access. For each resource, 
copy the stable Canvasflow resource identifier displayed in the 
UI (format: \texttt{cf:<type>:<id>}, e.g., \texttt{cf:issue:123}). These 
identifiers are used to build the mapping table in the IdP hook (Step 2.3).

Example of what the dashboard exposes:

\begin{lstlisting}
Publication: "Quarterly Review"  ->  cf:issue:123
Publication: "Weekly Brief"      ->  cf:issue:201
\end{lstlisting}

Identifiers are stable and will not change (\hyperref[sec:canvasflow-guarantees]{\S11 Canvasflow Guarantees~6}). 
They may be stored in the tenant's own systems for future reference.

\textbf{Step~2.2 --- Register the Canvasflow SPA as an OAuth client.}

Configure the client registration in the tenant's IdP with the following mandatory
settings:

\begin{table}[htbp]
\centering
\begin{tabularx}{\textwidth}{@{}lX@{}}
\toprule
\textbf{Setting} & \textbf{Required value} \\
\midrule
Grant type & Authorization Code \\
PKCE & Required, S256 method \\
Redirect URI & The tenant-specific callback URI:\texttt{https://www.<tenant-domain>/callback} — use the tenant's actual domain \\
Allowed scopes & \texttt{openid} at minimum \\
Token endpoint authentication & None (public client --- no client secret) \\
Access token format & JWT with \texttt{typ: "at+jwt"} in the header \\
Access token lifetime & $\leq$ 3600~seconds \\
\bottomrule
\end{tabularx}
\caption{Canvasflow SPA client registration settings}
\end{table}

Auth0 tenants MUST also register an API with \texttt{identifier = canvasflow-api}
(\hyperref[sec:api-registration-required]{\S6.3.4}).



\textbf{Step~2.3 --- Implement and test the pre-token hook.}



Implement the hook per the per-IdP guidance in \hyperref[sec:per-idp-implementation]{\S6.2} ,
following all four hook rules (\hyperref[sec:hook-rules]{\S6.4}). The hook must:
\begin{itemize}
  \item Inject \texttt{cf:entitlements} (\texttt{cf:entitlements} is RECOMMENDED 
  for all new integrations) into the
    \textbf{access token} (not the ID token)
  \item Re-evaluate entitlements from source on every login AND every refresh
  \item Emit an explicit empty array for users with no premium entitlements
  \item Produce values in \texttt{cf:<type>:<id>} format only, using the identifiers 
  copied from the Canvasflow administration dashboard in Step 2.1
\end{itemize}

After implementing, perform a test login and decode the resulting access token using
\url{https://jwt.io} or the Canvasflow OAuth Validator. Confirm that the entitlement
claim is present in the access token payload, not the ID token.

\textbf{Step~2.4 --- Configure signing and confirm credentials.}

\textit{JWKS path:} Confirm the IdP's JWKS URL. Verify that:
\begin{itemize}
  \item Issued tokens include a \texttt{kid} header matching a key in the JWKS
  \item The signing algorithm is RS256 or ES256
  \item RSA keys are at minimum 2048~bits
\end{itemize}

\textit{HS256 path (exception only):} Generate a cryptographically random secret of at
minimum 32~bytes. Confirm the \texttt{kid} header is present in issued tokens.

\textbf{Step~2.5 --- Share credentials with Canvasflow.}

Provide Canvasflow with:
\begin{itemize}
  \item \textbf{JWKS path:} the JWKS URL (public --- no secure channel required)
  \item \textbf{HS256 path:} the secret via a secure channel (encrypted message, shared
    secrets manager, or a call with screen share --- never unencrypted email)
  \item The exact \texttt{iss} value as it appears in issued tokens (trailing slash is
    significant)
  \item The \texttt{client\_id} registered in Step~2.2
  \item The tenant-specific callback URI registered in Step~2.2
    (e.g., \linebreak \texttt{https://www.tenant-domain.com/callback})
\end{itemize}

Canvasflow updates the \texttt{tenant\_config} row with these values and confirms
receipt.

\pagebreak

\subsection{Phase~3 --- Validation Using the Canvasflow OAuth Validator}

\subsubsection{The Validator Tool}

The Canvasflow OAuth Validator (\url{https://oauth.canvasflow.work}) is a browser-based
tool that executes the real Authorization Code + PKCE flow against the tenant's IdP ---
no simulation, no mocked responses, no OAuth libraries. It uses the Web Crypto API
directly for PKCE, decodes and annotates every JWT claim with Canvasflow-specific
checks, verifies signatures server-side (JWKS fetch or shared
HMAC), normalizes entitlements across the three claim names in precedence order, and
exercises the full refresh grant cycle.

The source code for the validator is available at
\url{https://github.com/Canvasflow/oauth-validator}.

\textbf{Why the validator is required.} The validator exercises the \emph{actual}
browser redirect and the \emph{actual} issued token --- not a simulated version. It
catches the entire class of \textit{works in jwt.io, fails in production} misconfigurations
before they affect a real user.

\begin{quote}
  \textit{Issues such as a claim placed in the ID token instead of the access token, 
  non-prefixed entitlement values, missing `typ`, or a broken refresh flow are 
  all detected and flagged by the validator with specific, 
  actionable error messages.}
\end{quote}

\subsubsection{Validation Checks}

The tenant's technical lead configures the validator with their IdP settings and MUST
produce all-green results on all five checks below. A partial pass --- any check showing
yellow (warning) or red (error) --- is not sufficient for go-live.

\textbf{V1 --- Token Structure.} The validator decodes the issued access token and
checks:
\begin{itemize}
  \item \texttt{typ: "at+jwt"} present in the header (or, for Cognito tenants:
    \texttt{typ: "JWT"} with \texttt{token\_use: "access"})
  \item \texttt{kid} present in the header
  \item \texttt{iss} exactly matches the registered issuer value
  \item \texttt{aud} contains \texttt{"canvasflow-api"}
  \item Entitlement claim present under one of the three supported names, as an array
  \item $\texttt{exp} - \texttt{iat} \leq 3600$~seconds
\end{itemize}

\textbf{V2 --- Signature Verification.} The validator's server-side component verifies
the token signature using the tenant's declared signing path (JWKS or HS256). The check
MUST pass without error.

\textbf{V3 --- Entitlement Normalization and Value Format.} The validator resolves the
entitlement claim using the three-claim precedence chain and applies the \texttt{cf:}
prefix filter. It confirms:
\begin{itemize}
  \item All expected Canvasflow resource IDs (from the Phase~1 mapping document) are
    present in the effective entitlement set
  \item Every value in the resolved claim carries the mandatory \texttt{cf:} prefix
  \item Any non-prefixed values are flagged as values that will be silently ignored in
    production (they do not block the check but are highlighted as 
    probable misconfigurations)
\end{itemize}

\begin{quote}
  \textit{A token whose effective entitlement set is empty — because all values 
  fail the `cf:` filter, or the claim is empty, or the claim is absent — 
  grants access to free content only. This is the correct baseline behavior 
  and is not a V3 failure condition, though the validator highlights missing 
  claims and non-prefixed values as probable misconfigurations the tenant 
  SHOULD resolve.}
\end{quote}

\textbf{V4 --- Refresh Flow.} The validator performs a refresh grant using the refresh
token obtained in V1 and verifies:
\begin{itemize}
  \item The refresh grant succeeds (HTTP~200 from the token endpoint)
  \item The new access token passes V1 and V2 checks
  \item The new access token carries a freshly evaluated entitlement claim (not copied
    from the previous token)
\end{itemize}

\textbf{V5 --- Expiry Behavior.} The validator verifies that:
\begin{itemize}
  \item An expired access token produces a 401 response from the Canvasflow GraphQL API
  \item After the refresh flow, a new valid access token restores API access
\end{itemize}

\subsubsection{Evidence Submission}

After all five checks pass in the validator, the tenant's technical lead exports or
screenshots the validator results and submits them to Canvasflow as evidence. Canvasflow
performs a joint review of the validator evidence.

Only after confirming all five checks are green does Canvasflow set \linebreak
\texttt{tenant\_config.status = 'active'}. The tenant is notified of production go-live
upon this status change.

\textbf{No tenant goes live without passing the validator.} The validator gate is
non-negotiable. It protects both the tenant's users (from a broken integration that
degrades to a silent authentication failure) and Canvasflow's platform (from improperly
formatted tokens reaching production).
